% Ad M. Brutum 1.13 (ad-brutum-01-13) --- English translation
% Translated by Claude (Anthropic) from the Latin (Perseus).
% Style guide: STYLE.md.

\ciceroLetterOpener{Brutus Ciceroni salutem}

\ciceroSection{1} About M. Lepidus, the fear of the rest forces me to be afraid myself. If he should tear himself away from us --- which I should wish men have suspected of him rashly and unjustly --- I beg and beseech you, Cicero, calling to witness our bond of family and your goodwill towards me: forget that my sister's children are sons of Lepidus, and consider that I have succeeded to them in their father's place. If I obtain this from you, surely you will not hesitate to take up any task for them. Some men live with their relatives in one way, others in another; I, in the case of my sister's children, can do nothing by which either my own wish or my duty might be filled. What good can they confer on me --- if indeed we are worthy of any benefit being conferred --- or what shall I be able to render to my mother, my sister, and those boys, if no weight has counted with you and with the rest of the Senate, in opposition to their father Lepidus, on the part of their uncle Brutus?

\ciceroSection{2} I cannot, nor ought I, to write much to you, between my anxiety and my exasperation. For if, in a matter so grave and so binding, I have need of words to rouse you and steady you, then there is no hope you will do what I wish and what is right. Therefore look for no long entreaties; look at me myself, who am bound to obtain this from you --- whether from Cicero, as the man most closely tied to me, in private; or from a consular of such standing, family connection set aside, in his public capacity. What you mean to do, I should like you to write back to me as soon as possible. 1 July, from camp.
